% ============================================================
% Paper 89: Mersenne Primes, Even Perfect Numbers, and Infinitude
% Author: Michael Fuhrmann
% Specialist Mathematics Project, Build 168
% Date: 2026-03-12
% ============================================================
\documentclass[12pt,a4paper]{amsart}

\usepackage{amsmath,amssymb,amsthm,mathtools,enumitem,booktabs,array,hyperref}
\usepackage[utf8]{inputenc}
\usepackage[T1]{fontenc}
\usepackage{geometry}
\geometry{margin=2.5cm}

% ---- Theorem environments ----
\newtheorem{theorem}{Theorem}[section]
\newtheorem{lemma}[theorem]{Lemma}
\newtheorem{corollary}[theorem]{Corollary}
\newtheorem{proposition}[theorem]{Proposition}
\theoremstyle{definition}
\newtheorem{definition}[theorem]{Definition}
\newtheorem{example}[theorem]{Example}
\theoremstyle{remark}
\newtheorem{remark}[theorem]{Remark}
\newtheorem{conjecture}[theorem]{Conjecture}

% ---- Custom operators ----
\newcommand{\N}{\mathbb{N}}
DeclareMathOperator{\ord}{ord}
\newcommand{\Z}{\mathbb{Z}}
\newcommand{\Q}{\mathbb{Q}}
\newcommand{\R}{\mathbb{R}}
\newcommand{\F}{\mathbb{F}}

% ============================================================
\begin{document}

\title{Mersenne Primes, Even Perfect Numbers, and Infinitude}

\author{Michael Fuhrmann}
\address{Specialist Mathematics Project, Build 168}
\email{specialist-maths@project.local}
\date{2026-03-12}

\begin{abstract}
The Euclid--Euler theorem establishes a perfect bijection between even perfect
numbers and Mersenne primes: every even perfect number has the form
$2^{p-1}(2^p - 1)$ where $2^p - 1$ is prime.
As of October 2024, exactly 52 Mersenne primes are known, the largest being
$M_{136\,279\,841} = 2^{136\,279\,841} - 1$, discovered on October 12, 2024
by Luke Durant via the Great Internet Mersenne Prime Search (GIMPS) —
the first Mersenne prime with exponent exceeding $10^8$.
The central open question is whether infinitely many Mersenne primes
(equivalently, even perfect numbers) exist.
In this paper we present the Euclid--Euler theorem, the Lucas--Lehmer
primality test for Mersenne numbers, a survey of known Mersenne primes,
density heuristics due to Lenstra and Wagstaff, the connection to cyclotomic
polynomials and Wagstaff's repunit conjecture, and a discussion of the
infinitude problem.
\end{abstract}

\maketitle

\tableofcontents

% ============================================================
\section{Introduction}
% ============================================================

A \emph{Mersenne number} is an integer of the form $M_n = 2^n - 1$.
When $M_p = 2^p - 1$ is prime (requiring $p$ itself to be prime), it is
called a \emph{Mersenne prime}.
The first few are
\[
  M_2 = 3, \quad M_3 = 7, \quad M_5 = 31, \quad M_7 = 127, \quad
  M_{13} = 8191.
\]
Mersenne primes are intimately tied to perfect numbers through the following
classical theorem, proved in its full form by Euler:

\begin{theorem}[Euclid--Euler, complete]
\label{thm:euclideuler}
An even positive integer $n$ is perfect if and only if
$n = 2^{p-1}(2^p - 1)$ where $2^p - 1$ is a Mersenne prime.
\end{theorem}

This perfect correspondence --- one of the most elegant results in number
theory --- reduces the classification of even perfect numbers to the study
of Mersenne primes.
The question of whether there are infinitely many remains one of the most
celebrated open problems in mathematics.

\begin{conjecture}[Infinitely many Mersenne primes]
\label{conj:inf}
There are infinitely many Mersenne primes, hence infinitely many even
perfect numbers.
\end{conjecture}

% ============================================================
\section{The Euclid--Euler Theorem}
% ============================================================

\subsection{Euclid's direction}

\begin{theorem}[Euclid, Elements IX.36, $\approx$ 300 BCE]
If $2^p - 1$ is prime, then $n = 2^{p-1}(2^p - 1)$ is perfect.
\end{theorem}

\begin{proof}
Let $q = 2^p - 1$ be prime. Then
\[
  \sigma(n) = \sigma(2^{p-1}) \cdot \sigma(q)
  = (2^p - 1)(q + 1)
  = (2^p - 1) \cdot 2^p
  = 2 \cdot 2^{p-1}(2^p - 1) = 2n. \qedhere
\]
\end{proof}

\subsection{Euler's direction}

\begin{theorem}[Euler, $\approx$ 1747]
Every even perfect number has the form $2^{p-1}(2^p - 1)$ where $2^p - 1$
is prime.
\end{theorem}

\begin{proof}
Let $n$ be an even perfect number. Write $n = 2^{k-1} m$ where $m$ is odd
and $k \geq 2$.
Since $\sigma$ is multiplicative and $\gcd(2^{k-1}, m) = 1$:
\[
  \sigma(n) = \sigma(2^{k-1})\sigma(m) = (2^k - 1)\sigma(m).
\]
Setting $\sigma(n) = 2n = 2^k m$:
\[
  (2^k - 1)\sigma(m) = 2^k m.
\]
Let $d = \gcd(2^k - 1, 2^k) = 1$. Since $2^k - 1$ is odd and $2^k$ is a
power of $2$, they are coprime, so $(2^k - 1) \mid m$.
Write $m = (2^k - 1) \cdot t$. Substituting:
\[
  (2^k - 1)\sigma(m) = 2^k(2^k-1)t \implies \sigma(m) = 2^k t.
\]
But $m = (2^k-1)t$ and $\sigma(m) = 2^k t = m + t = (2^k - 1)t + t$.
This means $\sigma(m) = m + t$, so $m$ has exactly two divisors: $m$ and $t$
(assuming $t < m$).
For $m$ to have exactly two divisors it must be prime, requiring $t = 1$.
Thus $m = 2^k - 1$ is prime, and $n = 2^{k-1}(2^k - 1)$.
Setting $p = k$ completes the proof.
\end{proof}

\begin{corollary}
The even perfect numbers are in bijection with the Mersenne primes via
$M_p \mapsto 2^{p-1} M_p$.
\end{corollary}

% ============================================================
\section{A Necessary Condition: The Exponent Must Be Prime}
% ============================================================

\begin{lemma}
\label{lem:exprime}
If $M_n = 2^n - 1$ is prime, then $n$ is prime.
\end{lemma}

\begin{proof}
If $n = ab$ with $1 < a, b < n$, then
$2^n - 1 = (2^a)^b - 1 = (2^a - 1)((2^a)^{b-1} + \cdots + 1)$,
so $2^a - 1 \mid 2^n - 1$ and $1 < 2^a - 1 < 2^n - 1$, contradicting primality.
\end{proof}

\begin{remark}
The converse of Lemma~\ref{lem:exprime} is false: $M_{11} = 2047 = 23 \times 89$
is composite despite $11$ being prime.
\end{remark}

% ============================================================
\section{The Lucas--Lehmer Primality Test}
% ============================================================

The primary tool for testing whether $M_p = 2^p - 1$ is prime is the
Lucas--Lehmer test, which is far more efficient than general primality tests
for numbers of this special form.

\begin{definition}
Define the \emph{Lucas--Lehmer sequence} by
\[
  s_0 = 4, \qquad s_{k+1} = s_k^2 - 2 \pmod{M_p}.
\]
\end{definition}

\begin{theorem}[Lucas 1878, Lehmer 1930]
\label{thm:LL}
For an odd prime $p$, $M_p = 2^p - 1$ is prime if and only if
$s_{p-2} \equiv 0 \pmod{M_p}$.
\end{theorem}

\begin{proof}[Proof outline]
The proof relies on the theory of Lucas sequences.
Consider the ring $\Z[\omega]$ where $\omega = 2 + \sqrt{3}$, so that
$s_k = \omega^{2^k} + \omega^{-2^k}$ (working modulo $M_p$).
If $M_p$ is prime, by Fermat's little theorem arguments,
$\omega^{M_p+1} = \omega^{2^p} \equiv 1 \pmod{M_p}$
in the ring $\F_{M_p}[\sqrt{3}]$.
This leads to $s_{p-2} \equiv 0 \pmod{M_p}$.
The converse uses the order of $\omega$ modulo $M_p$.
\end{proof}

\begin{remark}
The Lucas--Lehmer test has time complexity $O(p^2 \log p \log \log p)$
(using fast multiplication), making it practical even for exponents $p$
in the tens of millions.
It is the test used by GIMPS.
\end{remark}

\begin{example}
Test $M_5 = 31$: sequence is $s_0 = 4$, $s_1 = 14$, $s_2 = 194 \equiv 194 - 6\cdot31 = 8$,
$s_3 = 64 - 2 = 62 \equiv 0 \pmod{31}$. Since $s_{5-2} = s_3 = 0$, $M_5 = 31$ is prime. $\checkmark$
\end{example}

% ============================================================
\section{Known Mersenne Primes}
% ============================================================

As of October 2024, exactly 52 Mersenne primes are known.
Their exponents $p$ are:

\begin{center}
\begin{tabular}{rlrl}
\toprule
\# & Exponent $p$ & \# & Exponent $p$ \\
\midrule
1 & 2 & 2 & 3 \\
3 & 5 & 4 & 7 \\
5 & 13 & 6 & 17 \\
7 & 19 & 8 & 31 \\
9 & 61 & 10 & 89 \\
11 & 107 & 12 & 127 \\
13 & 521 & 14 & 607 \\
15 & 1279 & 16 & 2203 \\
17 & 2281 & 18 & 3217 \\
19 & 4253 & 20 & 4423 \\
21 & 9689 & 22 & 9941 \\
23 & 11213 & 24 & 19937 \\
25 & 21701 & 26 & 23209 \\
27 & 44497 & 28 & 86243 \\
29 & 110503 & 30 & 132049 \\
31 & 216091 & 32 & 756839 \\
33 & 859433 & 34 & 1257787 \\
35 & 1398269 & 36 & 2976221 \\
37 & 3021377 & 38 & 6972593 \\
39 & 13466917 & 40 & 20996011 \\
41 & 24036583 & 42 & 25964951 \\
43 & 30402457 & 44 & 32582657 \\
45 & 37156667 & 46 & 42643801 \\
47 & 43112609 & 48 & 57885161 \\
49 & 74207281 & 50 & 77232917 \\
51 & 82589933 & 52 & 136279841 \\
\bottomrule
\end{tabular}
\end{center}

\begin{remark}
The 51st Mersenne prime, $M_{82\,589\,933}$, was discovered on December 7, 2018
by Patrick Laroche as part of GIMPS.
It has $24\,862\,048$ decimal digits.
The ordering of the last several Mersenne primes (especially 48--51) reflects
discovery order, not necessarily size order, as not all exponents have been
tested by GIMPS.
\end{remark}

\begin{remark}[52nd Mersenne Prime, October 2024]
\label{rem:52nd}
On October 12, 2024, GIMPS announced the discovery of the 52nd known Mersenne
prime:
\[
    M_{136\,279\,841} = 2^{136\,279\,841} - 1.
\]
This is the first Mersenne prime with exponent exceeding $10^8$.  It has
$41\,024\,320$ decimal digits, making it the largest known prime as of the
time of writing.  The discovery was made by Luke Durant using a GPU-accelerated
GIMPS client~\cite{GIMPS2024b}.
\end{remark}

% ============================================================
\section{The Great Internet Mersenne Prime Search (GIMPS)}
% ============================================================

\begin{definition}
GIMPS is a distributed computing project founded by George Woltman in 1996
that coordinates the testing of Mersenne numbers for primality using the
Lucas--Lehmer test.
\end{definition}

GIMPS uses the following strategy:
\begin{enumerate}[label=(\arabic*)]
  \item Assign untested exponents $p$ to volunteers' computers.
  \item Run the Lucas--Lehmer test (or trial division for small factors first).
  \item Use double-checking to verify results.
  \item Report verified primes.
\end{enumerate}

As of 2024, GIMPS has discovered 17 of the largest known Mersenne primes
(starting with $M_{1\,257\,787}$ in 1996).
The project won the Electronic Frontier Foundation's \$100\,000 prize for the
first 10-million-digit prime ($M_{43\,112\,609}$, 2008).
The \$150\,000 prize for the first 100-million-digit prime remains unclaimed
as of 2024: the largest known Mersenne prime $M_{136\,279\,841}$ has
approximately $41$ million digits.

% ============================================================
\section{Density Heuristics: The Lenstra--Wagstaff Conjecture}
% ============================================================

\begin{conjecture}[Lenstra--Wagstaff]
\label{conj:LW}
The number of Mersenne primes $M_p = 2^p - 1$ with $p \leq x$ is
asymptotically
\[
  \sim \frac{e^\gamma}{\log 2} \cdot \log\log x
  \approx 2.5695 \cdot \log\log x,
\]
where $\gamma \approx 0.5772$ is the Euler--Mascheroni constant.
\end{conjecture}

\begin{remark}
This heuristic arises from modelling primality of $M_p$ probabilistically.
By the prime number theorem, $M_p \approx 2^p$ has about $p \log 2$ digits,
so its ``probability'' of being prime is approximately $1/(p \log 2)$.
Summing over prime $p \leq x$:
\[
  \sum_{\substack{p \leq x \\ p \text{ prime}}} \frac{1}{p \log 2}
  \approx \frac{1}{\log 2} \log\log x.
\]
The factor $e^\gamma$ arises from a more careful analysis accounting for
the special structure of Mersenne numbers.
\end{remark}

\begin{remark}
Conjecture~\ref{conj:LW} implies infinitely many Mersenne primes (since
$\log\log x \to \infty$), consistent with Conjecture~\ref{conj:inf}.
However, the heuristic cannot substitute for a proof.
\end{remark}

% ============================================================
\section{Cyclotomic Polynomials and Mersenne Numbers}
% ============================================================

\begin{definition}
The \emph{$n$-th cyclotomic polynomial} is
\[
  \Phi_n(x) = \prod_{\substack{k=1 \\ \gcd(k,n)=1}}^n (x - e^{2\pi i k/n}).
\]
\end{definition}

Mersenne numbers satisfy $2^n - 1 = \prod_{d \mid n} \Phi_d(2)$.

\begin{proposition}
The factorisation
\[
  2^n - 1 = \prod_{d \mid n} \Phi_d(2)
\]
shows that any prime factor of $\Phi_p(2) = 2^{p-1} + \cdots + 1$ (for prime $p$)
must be $\equiv 1 \pmod{p}$.
\end{proposition}

\begin{proof}
A prime $q \mid \Phi_p(2)$ satisfies $\ord_q(2) = p$ (the order of $2$ mod $q$
divides $p$ but not $1$, so equals $p$). By Fermat, $p \mid q-1$, i.e.\
$q \equiv 1 \pmod{p}$.
\end{proof}

\begin{corollary}
Every prime factor of $M_p = 2^p - 1$ is $\equiv 1 \pmod{p}$
(for odd prime $p$).
This severely limits the candidates and is exploited in trial division.
\end{corollary}

% ============================================================
\section{Wagstaff's Repunit Connection}
% ============================================================

\begin{definition}
A \emph{Mersenne number base $b$} is $M_{p,b} = \frac{b^p - 1}{b - 1}$
for $b \geq 2$.
For $b = 2$: $M_{p,2} = 2^p - 1$.
For $b = 10$: the repunits $R_p = \underbrace{11\cdots1}_p$.
\end{definition}

Wagstaff studied the analogous problem for repunit primes (base 10 and other
bases) and formulated an extended density conjecture:

\begin{conjecture}[Wagstaff 1983]
The number of primes of the form $(b^p - 1)/(b-1)$ with prime $p \leq x$
is $\sim (e^\gamma / \log b) \log\log x$ for each fixed $b \geq 2$.
\end{conjecture}

\begin{remark}
This conjecture unifies the Mersenne prime heuristic across different bases
and further supports the expectation that infinitely many primes of each type
should exist.
\end{remark}

% ============================================================
\section{Properties of Even Perfect Numbers}
% ============================================================

\begin{proposition}
Every even perfect number $n = 2^{p-1}(2^p-1)$ satisfies:
\begin{enumerate}[label=(\roman*)]
  \item $n$ is triangular: $n = T_{2^p - 1} = \binom{2^p}{2}/1$, i.e.\
        $n = 1 + 2 + \cdots + (2^p - 1)$;
  \item $n$ is a sum of the first $2^{(p-1)/2}$ consecutive odd cubes
        (for odd $p$);
  \item The binary representation of $n$ consists of $p$ ones followed by
        $p-1$ zeros: $\underbrace{1\cdots1}_{p}\underbrace{0\cdots0}_{p-1}$;
  \item $n \equiv 1 \pmod{9}$ for $p > 3$;
  \item The digits of $n$ sum to $1 \pmod{9}$ (for $p > 3$).
\end{enumerate}
\end{proposition}

\begin{proof}[Proof of (i)]
$\sum_{k=1}^{2^p - 1} k = \frac{(2^p-1)2^p}{2} = 2^{p-1}(2^p - 1) = n$.
\end{proof}

\begin{proof}[Proof of (iii)]
In binary: $2^{p-1}(2^p - 1) = 2^{2p-1} - 2^{p-1}$. Since $2^p - 1 = 2^{p-1} + \cdots + 2^0$
in binary is $p$ ones, multiplying by $2^{p-1}$ shifts left $p-1$ places, giving
$p$ ones starting at bit $p-1$.
\end{proof}

% ============================================================
\section{Gaps Between Mersenne Primes and Analytic Aspects}
% ============================================================

\begin{remark}
The known Mersenne prime exponents grow roughly exponentially.
If the Lenstra--Wagstaff conjecture is correct, one expects approximately one
new Mersenne prime per 1.5-fold increase in the search depth.
\end{remark}

\begin{proposition}[Carmichael's theorem for Mersenne numbers]
For $n > 2$, $M_n = 2^n - 1$ has a prime factor that does not divide
$M_k$ for any $1 \leq k < n$ (a \emph{primitive prime divisor}).
\end{proposition}

\begin{proof}
This follows from Zsygmondy's (Birkhoff--Vandiver) theorem applied to
$a = 2$, $b = 1$, $n \geq 3$.
\end{proof}

\begin{theorem}[Distribution of Mersenne prime exponents]
Under the Lenstra--Wagstaff heuristic, the expected number of Mersenne prime
exponents in the interval $[x, x + x/\log x]$ tends to $e^\gamma / \log 2$
as $x \to \infty$.
\end{theorem}

% ============================================================
\section{Connections to Other Structures}
% ============================================================

\subsection{Connection to $2$-perfect numbers}

Even perfect numbers satisfy $\sigma(n) = 2n$ and are linked to the study of
$2$-perfect (doubly perfect) numbers.
The Euclid--Euler bijection is the only known complete classification of a
class of perfect numbers.

\subsection{Connection to amicable numbers}

Two numbers $m, n$ are \emph{amicable} if $s(m) = n$ and $s(n) = m$.
Euler found 58 amicable pairs and conjectured infinitely many; this too is open.

\subsection{Wieferich primes}

A prime $p$ is a \emph{Wieferich prime} if $2^{p-1} \equiv 1 \pmod{p^2}$.
Only two are known: $1093$ and $3511$.
The connection to Mersenne primes: if $M_p$ is prime, then $p$ is typically
not Wieferich (heuristically), though no proof is known.

% ============================================================
\section{The Infinitude Problem}
% ============================================================

\begin{conjecture}[Restated: Infinitely many Mersenne primes]
There are infinitely many primes of the form $2^p - 1$.
\end{conjecture}

Despite extensive heuristic evidence and computational support, no proof
of infinitude is known.
The key difficulty is that the primes $p$ in $M_p = 2^p - 1$ must simultaneously
satisfy:
\begin{itemize}[nosep]
  \item $p$ itself is prime;
  \item $M_p$ passes the Lucas--Lehmer test.
\end{itemize}

\begin{remark}[Why it is hard]
The infinitude of primes of the form $f(n)$ for most polynomials $f$ (even
linear ones, i.e.\ Dirichlet's theorem) requires deep analytic methods.
For exponential functions like $2^p - 1$, current analytic number theory has
no tools to prove infinitude directly.
Sieve methods are ineffective because the sequence grows too fast.
\end{remark}

\begin{remark}[Upper bound impossibility]
It is also unknown whether the set of Mersenne primes is \emph{finite}.
Proving the set is finite would require showing that the Lucas--Lehmer test
eventually always fails, which seems equally intractable.
\end{remark}

% ============================================================
\section{Conclusion}
% ============================================================

The study of Mersenne primes and even perfect numbers sits at the intersection
of ancient number theory (Euclid's construction) and cutting-edge computational
mathematics (GIMPS, Lucas--Lehmer).
The Euclid--Euler theorem provides a complete and elegant classification of
even perfect numbers, reducing the problem to the study of primes of the
form $2^p - 1$.

The central questions remain:
\begin{enumerate}[label=(\arabic*)]
  \item Are there infinitely many Mersenne primes?
  \item Are there any odd perfect numbers? (see Paper~88)
  \item Can the Lenstra--Wagstaff density conjecture be proved?
  \item Is there a pattern among the exponents of Mersenne primes?
\end{enumerate}

These questions, seemingly simple to state, connect to some of the deepest
unsolved problems in analytic and computational number theory.

% ============================================================
\begin{thebibliography}{99}

\bibitem{Euclid300}
Euclid,
\textit{Elements},
Book IX, Proposition 36, $\approx$ 300 BCE.

\bibitem{Euler1747}
L.\ Euler,
\textit{De numeris amicabilibus},
Nova Acta Eruditorum (1747).

\bibitem{Lucas1878}
E.\ Lucas,
\textit{Théorie des fonctions numériques simplement périodiques},
Amer.\ J.\ Math.\ \textbf{1} (1878), 184--196, 197--240, 289--321.

\bibitem{Lehmer1930}
D.\ H.\ Lehmer,
\textit{An extended theory of Lucas' functions},
Ann.\ Math.\ \textbf{31} (1930), 419--448.

\bibitem{Wagstaff1983}
S.\ S.\ Wagstaff Jr.,
\textit{Divisors of Mersenne numbers},
Math.\ Comp.\ \textbf{40} (1983), 385--397.

\bibitem{LenstraWagstaff}
H.\ W.\ Lenstra Jr.\ and R.\ D.\ Wagstaff Jr.,
\textit{Unpublished heuristic analysis cited in}:
P.\ Ribenboim, \textit{The New Book of Prime Number Records},
Springer, 1996.

\bibitem{Caldwell2024}
C.\ K.\ Caldwell,
\textit{Mersenne Primes: History, Theorems and Lists},
The Prime Pages, \url{https://primes.utm.edu/mersenne/}, 2024.

\bibitem{GIMPS2024}
GIMPS,
\textit{Great Internet Mersenne Prime Search},
\url{https://www.mersenne.org/}, 2024.

\bibitem{Laroche2018}
P.\ Laroche,
\textit{Announcement of $M_{82\,589\,933}$},
GIMPS Press Release, December 2018.

\bibitem{Zsygmondy1892}
K.\ Zsygmondy,
\textit{Zur Theorie der Potenzreste},
Monatsh.\ Math.\ Phys.\ \textbf{3} (1892), 265--284.

\bibitem{Crandall2005}
R.\ Crandall and C.\ Pomerance,
\textit{Prime Numbers: A Computational Perspective},
2nd ed., Springer, 2005.

\bibitem{GIMPS2024b}
GIMPS,
\textit{GIMPS Project Discovers Largest Known Prime Number: $2^{136\,279\,841}-1$},
Press release, October 12, 2024,
\url{https://www.mersenne.org/}.

\bibitem{Ribenboim1996}
P.\ Ribenboim,
\textit{The New Book of Prime Number Records},
Springer, New York, 1996.

\bibitem{Hardy2008}
G.\ H.\ Hardy and E.\ M.\ Wright,
\textit{An Introduction to the Theory of Numbers},
6th ed., Oxford University Press, 2008.

\end{thebibliography}

\end{document}
